% GENERATED FILE. Edit canonical lessons, then run npm run generate:books.
% canonical-source-hash: fnv1a64:622df1c555d5b4cc
% canonical-lessons: ZH-C115-xiaoxin, ZH-W115-xin-heart, ZH-C115-xiyan, ZH-C115-jinzhi, ZH-C115-zhuyi, ZH-C115-weixian

\chapter{Careful!}
\label{ch:zh-careful}
\hlchaptermodality{115}

\begin{chapteropening}
\textbf{By the end of this chapter:} \emph{I can say be careful, no smoking, pay attention and dangerous.}

Complete the last lesson of chapter 115: I can say be careful, no smoking, pay attention and dangerous.
\end{chapteropening}

\section[\texorpdfstring{\zh{小心} (xiǎoxīn)}{小心 (xiǎoxīn)}]{\zh{小心} (xiǎoxīn) — to be careful; to mind (something)}

\label{lesson:ZH-C115-xiaoxin}



\begin{quote}
Before the new one: say the Mandarin for to pay by card, then the Mandarin for to give change.
\end{quote}

\subsection*{You'll want to know: \zh{小心}}
\textbf{\zh{小心}} — \emph{xiǎoxīn} — \textquotedblleft{}to be careful; to mind (something)\textquotedblright{}.

\begin{quote}
\textbf{\zh{出门小心车}} \emph{chūmén xiǎoxīn chē} — Mind the traffic when you go out.
\end{quote}

The new character, in pieces: \zh{心} — four strokes: a dot, a long hook lying along the bottom, and two dots, the foot of \zh{想}. The next lesson puts it in your hand.

\begin{cousinweb}
\textbf{\zh{小心}} (\emph{xiǎoxīn}) is built from \zh{小} 'small' and \zh{心} 'heart', both inherited from Middle Chinese, \emph{sjewX} and \emph{sim} (Guangyun, in Baxter's transcription). How a small heart came to mean care is not settled.
\end{cousinweb}

\begin{grammarlens}[title={What you've built}]
One more word for signs and notices.
\end{grammarlens}

\subsection*{Your turn}
\begin{itemize}
\raggedright
  \item \emph{Say it:} \emph{xiǎoxīn}
  \item \emph{Say it:} \emph{xiǎoxīn}, once more
  \item \emph{Say it:} \emph{chūmén xiǎoxīn chē} — Mind the traffic when you go out.
  \item \emph{Recall:} say the Mandarin for to pay by card, then the Mandarin for to give change, then say \emph{xiǎoxīn} again
\end{itemize}

\begin{tcolorbox}[breakable,colback=teal!4,colframe=teal!35!black,title={Before you move on}]
What does \zh{小心} mean? (\textquotedblleft{}To be careful; to mind (something)\textquotedblright{}.) Say it once more.
\end{tcolorbox}
\section[\texorpdfstring{\zh{心} (xīn)}{心 (xīn)}]{\zh{心} — a dot, a lying hook, two dots}

\label{lesson:ZH-W115-xin-heart}



\begin{quote}
The word first, then the character inside it.

\begin{itemize}
\raggedright
  \item \emph{Recall:} say \emph{xiǎoxīn} — \textbf{R1}, one lesson back, before it decays
\end{itemize}

\begin{quote}
\zh{想} \zh{小}
\end{quote}

\emph{Read it:} each one aloud, with nothing said first — \emph{xiǎng}, \emph{xiǎo}
\end{quote}

\subsection*{Script — a dot, a lying hook, two dots}
\hlblockfigure{\detokenize{figures/ZH-W115-xin-heart-filmstrip.pdf}}{How \zh{心} is written, stroke by stroke}

\begin{quote}
\zh{心}
\end{quote}

\textbf{A dot} on the left, slanting down.

\textbf{A long hook} lying along the bottom, curving up at the right and hooking in.

\textbf{Two dots} above it, the middle one, then the right one.

1 + 1 + 2 = four.

\subsection*{Your turn}
\begin{itemize}
\raggedright
  \item \emph{Write it:} \zh{心}, all four, with the model in front of you
  \item \emph{Cover:} write it again from memory, and count the pieces, not the strokes
\end{itemize}

\begin{tcolorbox}[breakable,colback=teal!4,colframe=teal!35!black,title={Before you move on}]
How many strokes in \zh{心}? (\textbf{Four}.) Which stroke is the long one? (\textbf{The hook lying along the bottom, second}.) Where have you written it before? (\textbf{At the foot of \zh{想}}.)

Source: \href{https://raw.githubusercontent.com/chanind/hanzi-writer-data/68d10a4b21150cae5e1ebbd223eed289cf32d90c/data/\%E5\%BF\%83.json}{Hanzi Writer Data \zh{心}, pinned stroke paths and medians}.
\end{tcolorbox}
\section[xīyān]{xīyān — to smoke}

\label{lesson:ZH-C115-xiyan}



\begin{quote}
Before the new one: say the Mandarin for to give change, then the Mandarin for to be careful.
\end{quote}

\subsection*{You'll want to know: xīyān}
\emph{xīyān} — \textquotedblleft{}to smoke\textquotedblright{}.

\begin{quote}
\emph{nǐ xīyān ma} — Do you smoke?
\end{quote}

Its characters come later; for now it is a word for the ear.

\begin{cousinweb}
\emph{xīyān} is built from xī 'to breathe in' and yān 'smoke', both inherited from Middle Chinese, \emph{xip} and \emph{'en} (Guangyun, in Baxter's transcription).
\end{cousinweb}

\begin{grammarlens}[title={What you've built}]
One more word for signs and notices.
\end{grammarlens}

\subsection*{Your turn}
\begin{itemize}
\raggedright
  \item \emph{Say it:} \emph{xīyān}
  \item \emph{Say it:} \emph{xīyān}, once more
  \item \emph{Say it:} \emph{nǐ xīyān ma} — Do you smoke?
  \item \emph{Recall:} say the Mandarin for to give change, then the Mandarin for to be careful, then say \emph{xīyān} again
\end{itemize}

\begin{tcolorbox}[breakable,colback=teal!4,colframe=teal!35!black,title={Before you move on}]
What does xīyān mean? (\textquotedblleft{}To smoke\textquotedblright{}.) Say it once more.
\end{tcolorbox}
\section[jìnzhǐ]{jìnzhǐ — to forbid; no (on signs)}

\label{lesson:ZH-C115-jinzhi}



\begin{quote}
Before the new one: say the Mandarin for to be careful, then the Mandarin for to smoke.
\end{quote}

\subsection*{You'll want to know: jìnzhǐ}
\emph{jìnzhǐ} — \textquotedblleft{}to forbid; no (on signs)\textquotedblright{}.

\begin{quote}
\emph{zhèr jìnzhǐ xīyān} — No smoking here.
\end{quote}

Its characters come later; for now it is a word for the ear.

\begin{cousinweb}
\emph{jìnzhǐ} is built from jìn and zhǐ, both inherited from Middle Chinese, \emph{kimH} and \emph{tsyiX} (Guangyun, in Baxter's transcription), and the dictionary glosses both 'to stop'.
\end{cousinweb}

\begin{grammarlens}[title={What you've built}]
One more word for signs and notices.
\end{grammarlens}

\subsection*{Your turn}
\begin{itemize}
\raggedright
  \item \emph{Say it:} \emph{jìnzhǐ}
  \item \emph{Say it:} \emph{jìnzhǐ}, once more
  \item \emph{Say it:} \emph{zhèr jìnzhǐ xīyān} — No smoking here.
  \item \emph{Recall:} say the Mandarin for to be careful, then the Mandarin for to smoke, then say \emph{jìnzhǐ} again
\end{itemize}

\begin{tcolorbox}[breakable,colback=teal!4,colframe=teal!35!black,title={Before you move on}]
What does jìnzhǐ mean? (\textquotedblleft{}To forbid; no (on signs)\textquotedblright{}.) Say it once more.
\end{tcolorbox}
\section[zhùyì]{zhùyì — to pay attention; to take care}

\label{lesson:ZH-C115-zhuyi}



\begin{quote}
Before the new one: say the Mandarin for to smoke, then the Mandarin for to forbid.
\end{quote}

\subsection*{You'll want to know: zhùyì}
\emph{zhùyì} — \textquotedblleft{}to pay attention; to take care\textquotedblright{}.

\begin{quote}
\emph{qǐng zhùyì nǐ de bāo} — Please keep an eye on your bag.
\end{quote}

Its characters come later; for now it is a word for the ear.

\begin{cousinweb}
\emph{zhùyì} is built from zhù 'to pour into' and yì 'mind', both inherited from Middle Chinese, \emph{tsyuH} and \emph{'iH} (Guangyun, in Baxter's transcription): to pour the mind into something.
\end{cousinweb}

\begin{grammarlens}[title={What you've built}]
One more word for signs and notices.
\end{grammarlens}

\subsection*{Your turn}
\begin{itemize}
\raggedright
  \item \emph{Say it:} \emph{zhùyì}
  \item \emph{Say it:} \emph{zhùyì}, once more
  \item \emph{Say it:} \emph{qǐng zhùyì nǐ de bāo} — Please keep an eye on your bag.
  \item \emph{Recall:} say the Mandarin for to smoke, then the Mandarin for to forbid, then say \emph{zhùyì} again
\end{itemize}

\begin{tcolorbox}[breakable,colback=teal!4,colframe=teal!35!black,title={Before you move on}]
What does zhùyì mean? (\textquotedblleft{}To pay attention; to take care\textquotedblright{}.) Say it once more.
\end{tcolorbox}
\section[wēixiǎn]{wēixiǎn — dangerous}

\label{lesson:ZH-C115-weixian}



\begin{quote}
Before the new one: say the Mandarin for to forbid, then the Mandarin for to pay attention.
\end{quote}

\subsection*{You'll want to know: wēixiǎn}
\emph{wēixiǎn} — \textquotedblleft{}dangerous\textquotedblright{}.

\begin{quote}
\emph{zhèr hěn wēixiǎn} — It's dangerous here.
\end{quote}

Its characters come later; for now it is a word for the ear.

\begin{cousinweb}
\emph{wēixiǎn} is built from wēi and xiǎn, both inherited from Middle Chinese, \emph{ngjwe} and \emph{xjemX} (Guangyun, in Baxter's transcription), glossed 'unsafe' and 'perilous, hard to pass'. The regular change from \emph{ngjwe} gives wéi, with a rising tone; the level wēi of today is not settled.
\end{cousinweb}

\begin{grammarlens}[title={What you've built}]
One more word for signs and notices.
\end{grammarlens}

\subsection*{Your turn}
\begin{itemize}
\raggedright
  \item \emph{Say it:} \emph{wēixiǎn}
  \item \emph{Say it:} \emph{wēixiǎn}, once more
  \item \emph{Say it:} \emph{zhèr hěn wēixiǎn} — It's dangerous here.
  \item \emph{Recall:} say the Mandarin for to forbid, then the Mandarin for to pay attention, then say \emph{wēixiǎn} again
\end{itemize}

\begin{tcolorbox}[breakable,colback=teal!4,colframe=teal!35!black,title={Before you move on}]
What does wēixiǎn mean? (\textquotedblleft{}Dangerous\textquotedblright{}.) Say it once more.
\end{tcolorbox}
